% problem-000065 7 荧光探针识别甲醛
\section*{7. 荧光探针识别甲醛}
\textit{下列为分问。}

\subsection*{7-1}
某些物质吸收光后，电⼦从基态跃迁⾄激发态，⽽由激发跃迁回基态时，以产⽣荧光的形式将能量释放出。研究发现， $\lambda _ { \operatorname* { m a x } } = 6 4 5 \mathrm { n m }$ 的光激发下，化合物1⼏乎不产⽣荧光，但是化合物1与甲醛反应⽣成的化合物2能产⽣很强的荧光。据此，化合物1能够⾼选择性识别甲醛。

（图：）

\subsection*{7-1}
-1 简要解释化合物1⼏乎不产⽣荧光，⽽化合物2能产⽣很强的荧光。

\subsection*{7-1}
-2 写出化合物1与甲醛反应⽣成化合物2的反应机理。

（图：）

（图：）

\subsection*{7-2}
化合物1和化合物2分别与(CH_{3})_{3}Al按摩尔⽐1:3投料，发⽣下列反应：

\subsection*{7-2}
-1 写出产物A，B的结构简式，标注B中不对称碳原⼦的绝对构型（�或�）。

\subsection*{7-2}
-2 写出⽣成A和B的反应机理，解释反应的区域选择性和⽴体选择性。

\subsection*{7-2}
-3 ⽤正⼰烷作溶剂（其他条件相同），A和B的产率很低，并有⼤量原料未反应。解释溶剂对反应过程的影响。
